\documentclass[11pt]{article}
\usepackage[margin=1in]{geometry}
\usepackage[T1]{fontenc}\usepackage[utf8]{inputenc}
\usepackage{lmodern,amsmath,amssymb,amsthm,mathtools,microtype}
\usepackage[colorlinks=true,linkcolor=blue,citecolor=blue,urlcolor=blue]{hyperref}
\hypersetup{pdftitle={Path-norm geometry of Gaussian ReLU posteriors},pdfauthor={ChatGPT-6 Astra; Claude Opus 5.5; Kaarel Hänni; Dmitry Vaintrob}}
\newtheorem{theorem}{Theorem}[section]
\newtheorem{lemma}[theorem]{Lemma}
\newtheorem{proposition}[theorem]{Proposition}
\newtheorem{corollary}[theorem]{Corollary}
\theoremstyle{remark}\newtheorem{remark}[theorem]{Remark}
\title{Path-norm geometry of Gaussian ReLU posteriors}
\author{ChatGPT-6 Astra \and Claude Opus 5.5 \and Kaarel H\"anni \and
Dmitry Vaintrob\thanks{Authors are listed in alphabetical order.}}
\date{17 November 2026}
\begin{document}\maketitle
\begin{abstract}
We study exact Bayesian regression with a one-hidden-layer ReLU network,
independent Gaussian weights, a Gaussian likelihood, and a raw linear
readout. For one-dimensional inputs, normalizing each incoming row
transforms the function prior exactly into independent uniform ridge
directions with independent Laplace amplitudes. This identifies a
quantitative connection to convex path-norm regression: when learner
width $N$ is at least sample size $n$, the actual posterior's expected empirical
loss plus path penalty exceeds the full convex optimum by at most
$C\lambda\nu\log(e/(\lambda\nu))$. Here $\nu\in(0,1]$ is the hidden
Gaussian variance and $\lambda\in(0,1/2]$ is the normalized likelihood
penalty. The same bound controls expected squared training-vector distance to
the unique convex fitted vector. A direct Gaussian argument extends
these empirical conclusions to every input dimension, with an explicit
dimension factor. Conditional on directions and training
predictions, excess amplitude mass has a Gamma-type tail. These facts
do not imply bounded predictions between observations. We construct
four bounded data pairs for which every minimum full-path-norm
interpolant has an arbitrarily large excursion; at an explicit
likelihood schedule, the ordinary Gaussian posterior mean does too,
uniformly over all eligible widths and scales. A final conditional
variance bound separates last-layer uncertainty from the unresolved
variation in hidden-feature posterior means. The results concern
finite architectures and do not change the prior or introduce an
architecture mixture.
\end{abstract}

\section{Empirical fit and predictions between observations}
A path penalty is often used to understand why an overparameterized
ReLU network can favor a simpler fit. For Bayesian learning, however,
three objects must be distinguished: the posterior's parameter
representation, the minimum path mass needed for its training
predictions, and its function away from the training inputs. The
results here quantify the first two relationships and give a precise
limitation of an inference from them to the third.

We use the ordinary posterior after the given sample. The input domain
is $[-1,1]$, and the regression output is unbounded: neither the
likelihood nor prediction uses an output clipping operation. All
results are deterministic in the observed sample except for the
posterior randomness being averaged. In particular, the range
counterexample does not assert that its four-point configuration is
likely under a well-spread IID design.

A normalized ridge is $x\mapsto(\theta_1x+\theta_2)_+$ with
$\theta_1^2+\theta_2^2=1$. The mass of an implementation is the sum
of absolute readout amplitudes. Minimizing that mass over all
implementations of a function gives its full path norm; minimizing
only subject to a training vector can give a different value. We
keep that distinction explicit in the statements below.

\section{Two posterior-specific reductions}
We first quantify how the posterior relates to convex path-norm
regression on the observed sample.
Let
\[
 z(x)=\sum_{j=1}^N v_j(a_jx+b_j)_+,
 \quad (a_j,b_j)\sim N(0,\nu I_2),\quad
 v_j\sim N(0,\nu/N^2),\quad \omega=N/\nu,
\]
where all coordinates are independent and $0<\nu\le1$.
The actual posterior has density proportional to
$\exp[-\omega\widehat R(z)/(2\lambda)]$ relative to this prior.
Throughout this section $|X_i|,|Y_i|\le1$.

Write $(a_j,b_j)=r_j\theta_j$ with $\theta_j$ uniform on the unit
circle, and put $c_j=v_jr_j$. The exact pushforward prior has independent
uniform directions and independent Laplace amplitudes
$c_j\sim\operatorname{Laplace}(0,\omega^{-1})$.
Indeed $r_j^2$ is exponential of mean $2\nu$, and
\[
 \mathbb E e^{itc_j}
 =\mathbb E\exp[-\nu t^2r_j^2/(2N^2)]
 =(1+\nu^2t^2/N^2)^{-1}.
\]
Put $P(c)=\sum_j|c_j|$ and
$A_\theta c=(\sum_jc_j(\theta_{j,1}X_i+\theta_{j,2})_+)_{i=1}^n$.
All statements below are about this exact pushforward posterior;
no replacement prior or estimator is introduced.

\begin{proposition}[Expected excess mass on a training-prediction fiber]
Fix the directions and a feasible training vector $b$.
Let $d_\theta=N-\operatorname{rank}(A_\theta)$ and
\[
 P_{\min}(b,\theta)=\min\{\|c\|_1:A_\theta c=b\}.
\]
For almost every fiber under the regular conditional posterior,
\[
 \mathbb E[P(c)-P_{\min}(b,\theta)\mid\theta,A_\theta c=b]
 \le d_\theta/\omega\le\nu. \tag{PR1}
\]
\end{proposition}
\begin{proof}
Conditional on the fiber, the likelihood is constant. The posterior
density with respect to its affine Hausdorff measure is proportional
to $e^{-\omega\|c\|_1}$; the coarea normalization is constant there.
Choose a minimizer $c_*$. In orthonormal nullspace coordinates $u$,
let $V(u)=\|c_*+Bu\|_1$, where $B$ is an isometry from
$\mathbb R^{d_\theta}$ onto $\ker A_\theta$.
The density is coercive and has exponential tails. Integration of
$\operatorname{div}(u e^{-\omega V(u)})$ gives
\[
 d_\theta=\omega\mathbb E[u\cdot\nabla V(u)].
\]
The piecewise-linear convex potential is differentiable almost
everywhere; the identity also follows by smooth approximation.
Convexity gives $u\cdot\nabla V(u)\ge V(u)-V(0)$.
The zero-dimensional case is immediate. This proves (PR1).
\end{proof}

\begin{proposition}[Approximate minimization of the full empirical path objective]
Suppose $N\ge n$ and $0<\lambda\le1/2$. Let $J_*$ be the minimum,
over all finite shallow ReLU implementations, of
\[
 J(c,\theta)=\widehat R(z_{c,\theta})+2\lambda\sum_j|c_j|.
\]
Equivalently, one may first minimize the amplitude mass among
implementations of each training vector, then minimize over those vectors. Then the actual posterior satisfies
\[
 0\le\mathbb E_\Pi[\widehat R(z)+2\lambda P(c)]-J_*
 \le C\lambda\nu\log\frac{e}{\lambda\nu}. \tag{PR2}
\]
Here $C$ is universal. The mass in the posterior expectation is that
of its actual amplitude representation, not a separately minimized
version.
\end{proposition}
\begin{proof}
The signed unit-ridge training vectors form a compact subset of
$\mathbb R^n$. The minimum-mass representation of any feasible
training vector therefore exists, and a linear-dependence reduction
uses at most $n$ atoms. Minimizing the coercive objective over the
resulting compact sublevel sets gives a minimizer $f_*$ with at most
$n$ atoms. Since zero is feasible,
\[
 J_*\le1,\qquad \widehat R(f_*)\le1,\qquad P_*\le(2\lambda)^{-1}.
\]
Thus $f_*$ fits into the prescribed width $N$.

Let $F=-(2\lambda/\omega)\log Z$ be the actual free energy.
Represent $f_*$ by the balanced Gaussian center
$\bar r_j=\sqrt{N|c_j|}\theta_j$,
$\bar v_j=\operatorname{sign}(c_j)\sqrt{|c_j|/N}$.
Its weighted Gaussian energy is $2P_*$.
Around this center use native covariance multiplied by
$\epsilon^2$, with $\epsilon=(\lambda\nu)^2$.
The comparison law has
\[
 \operatorname{KL}(\rho\Vert\mu)
 \le\omega P_*+3N\log(1/\epsilon).
\]
Its uniform root-mean-square function perturbation is bounded by
$C\epsilon\sqrt{\nu P_*}+C\epsilon^2\nu$.
Using $\widehat R(f_*)\le1$ gives
$\rho\widehat R\le\widehat R(f_*)+C\lambda\nu$.
The Gibbs variational inequality therefore yields
\[
 F\le J_*+C\lambda\nu\log\frac e{\lambda\nu}. \tag{PR3}
\]

For the pushforward prior,
$\mathbb E e^{t\omega P}=(1-t)^{-N}$ for $0<t<1$.
The entropy inequality and the Gibbs identity imply
\[
 F\ge\mathbb E_\Pi\widehat R+2\lambda t\mathbb E_\Pi P
       +2\lambda\nu\log(1-t).
\]
Radial integration by parts in the Laplace amplitudes gives
\[
 \mathbb E_\Pi P
 =\nu-\lambda^{-1}\mathbb E_\Pi P_n[z(z-Y)]
 \le\nu+\widehat{Y^2}/(4\lambda)
 \le\nu+(4\lambda)^{-1}.
\]
Use $1-t=\lambda\nu$. The lost coercivity term
$2\lambda(1-t)\mathbb E_\Pi P$ is at most $C\lambda\nu$.
Combining this with (PR3) proves the upper bound in (PR2).
Its lower bound follows from the definition of $J_*$.
\end{proof}

Neither statement alone controls population predictions in an empty
sample interval. In (PR1), a small mass excess need not imply small
function distance to a minimum-mass representative: nearby ridge
directions can make the relevant stability constant arbitrarily bad.
In (PR2), near-minimal objective value is not a uniform output-range
bound. The four-point example later in this paper proves that
full Gaussian path-norm interpolants themselves can have arbitrarily
large excursions between bounded training labels. A statistical
improvement must use the local sample spacings and the cost of fitting
nearby conflicting labels, rather than invoking deterministic range
preservation.

\begin{corollary}[Conditional excess-mass tails]
With the notation of (PR1), for every $0<s<1$,
\[
 \mathbb E\left[e^{s\omega(P-P_{\min})}\mid\theta,A_\theta c=b\right]
 \le(1-s)^{-d_\theta}.
\]
Consequently, for every $t>0$, the conditional probability that
\[
 P-P_{\min}>\omega^{-1}\{d_\theta+2\sqrt{d_\theta t}+2t\}
\]
is at most $e^{-t}$.
\end{corollary}
\begin{proof}
Write $W(u)=V(u)-V(0)$. It is convex, nonnegative, and vanishes at
zero. For $r=(1-s)^{-1}\ge1$, convexity gives $W(ru)\ge rW(u)$.
A change of variables therefore shows
\[
 \int e^{-\omega W(u)/r}\,du
 =r^{d_\theta}\int e^{-\omega W(ru)/r}\,du
 \le r^{d_\theta}\int e^{-\omega W(u)}\,du.
\]
This is the stated moment bound. The standard Chernoff calculation
for $(1-s)^{-d_\theta}$ gives the displayed, slightly relaxed,
Gamma-tail threshold. If $d_\theta=0$, the excess is identically zero.
\end{proof}

\begin{corollary}[Concentration around the convex fitted vector]
Under the assumptions of (PR2), let $b_*\in\mathbb R^n$ be the
training vector of a minimizer of the full empirical path objective.
This vector is unique, and
\[
 \mathbb E_\Pi\frac1n\|A_\theta c-b_*\|_2^2
 \le C\lambda\nu\log\frac e{\lambda\nu}. \tag{PR4}
\]
The same bound holds for the training vector of the ordinary
posterior mean.
\end{corollary}
\begin{proof}
Let $p(b)$ be the minimum signed-unit-ridge mass representing the
training vector $b$, allowing an arbitrary finite number of atoms.
Its domain is a linear subspace of $\mathbb R^n$ and $p$ is a norm
on that subspace. The objective
$\|b-Y\|_2^2/n+2\lambda p(b)$ is strictly convex there, so its
minimizing training vector is unique. The norm subgradient at its
minimum and the exact quadratic expansion give, for every feasible $b$,
\[
 \frac1n\|b-Y\|_2^2+2\lambda p(b)-J_*
 \ge\frac1n\|b-b_*\|_2^2.
\]
Every amplitude implementation has $P(c)\ge p(A_\theta c)$.
Apply (PR2) and then Jensen's inequality for the posterior mean.
\end{proof}

\section{Empirical convexification in any input dimension}
The objective and training-vector conclusions do not depend on the
one-dimensional Laplace identity. Gaussian entropy coercivity gives
the following extension directly. The fiber and range statements
remain specifically one-dimensional.

\begin{proposition}[A dimension-explicit empirical objective bound]
Let $d\ge1$, $x\in[-1,1]^d$, $\bar x=(x,1)$, and
\[
 z_w(x)=\sum_{j=1}^N v_j(w_j\cdot\bar x)_+,
 \qquad w_j\sim N(0,(\nu/d)I_{d+1}),\qquad
 v_j\sim N(0,\nu/N^2),
\]
independently. Put
\[
 Q=\frac dN\sum_j\|w_j\|_2^2+N\sum_jv_j^2,
 \quad P=\sqrt d\sum_j|v_j|\|w_j\|_2\le Q/2,
 \quad\omega=N/\nu.
\]
For $|Y_i|\le1$, $N\ge n$, $0<\nu\le1$, and $0<\lambda\le1/2$,
form the actual posterior proportional to
$e^{-\omega\widehat R(z)/(2\lambda)}d\mu$.
If $J_*$ denotes the minimum of empirical squared loss plus
$2\lambda$ times implementation mass, over all finite normalized-ridge
implementations, then
\[
 0\le\mathbb E_\Pi[\widehat R(z)+2\lambda P]-J_*
 \le C(d+2)\lambda\nu
          \log\frac{e(d+1)}{\lambda\nu}. \tag{GD1}
\]
The unique convex fitted training vector $b_*$ also satisfies
\[
 \mathbb E_\Pi\frac1n\|(z(X_i))_{i=1}^n-b_*\|_2^2
 \le C(d+2)\lambda\nu
          \log\frac{e(d+1)}{\lambda\nu}, \tag{GD2}
\]
and so does the training vector of the ordinary posterior mean.
\end{proposition}
\begin{proof}
Normalize a ridge as
$\phi_\theta(x)=d^{-1/2}(\theta\cdot\bar x)_+$ for
$\theta\in\mathbb S^d$. These features are bounded by $\sqrt2$,
and the signed amplitude is $c=\sqrt d\,v\|w\|_2$.
The compact signed-feature linear program again gives an optimum
$f_*$ with at most $n$ atoms and
$\widehat R(f_*)\le1$, $P_*\le(2\lambda)^{-1}$.
Balance an atom of amplitude $c$ with
\[
 \bar w=\sqrt{N|c|/d}\,\theta,
 \qquad \bar v=\operatorname{sign}(c)\sqrt{|c|/N}.
\]
The total center energy is $Q(\bar w)=2P_*$.

The parameter dimension is $D=(d+2)N$. A Gaussian comparison tube
with native covariance multiplied by $\epsilon^2$ has
\[
 \operatorname{KL}(\rho\Vert\mu)
 \le\omega P_*+D\log(1/\epsilon).
\]
Take $\epsilon=(\lambda\nu/(d+1))^2$.
The ReLU Lipschitz property, $\|\bar x\|_2^2\le d+1$, balance,
and independence of the incoming and readout perturbations give
uniform root-mean-square error at most
$C\epsilon\sqrt{\nu P_*}+C\epsilon^2\nu$.
The empirical loss increase is therefore at most $C\lambda\nu$.
The Gibbs variational principle gives
\[
 F\le J_*+C(d+2)\lambda\nu
                   \log\frac{e(d+1)}{\lambda\nu}.
\]

For $0<t<1$ the exact Gaussian moment is
\[
 \log\mathbb E_\mu e^{t\omega Q/2}
 =\frac D2\log\frac1{1-t}.
\]
The entropy inequality and $Q\ge2P$ imply
\[
 F\ge\mathbb E_\Pi\widehat R+2\lambda t\mathbb E_\Pi P
      -(d+2)\lambda\nu\log\frac1{1-t}.
\]
Gaussian radial integration by parts gives
\[
 \mathbb E_\Pi Q
 =(d+2)\nu-\frac2\lambda\mathbb E_\Pi P_n[z(z-Y)]
 \le(d+2)\nu+\frac1{2\lambda}.
\]
Use $1-t=\lambda\nu$ and $P\le Q/2$ to bound the lost coercivity
by $C(d+2)\lambda\nu$. This proves (GD1).
Strong convexity in the training vector proves (GD2) exactly as in
(PR4), and Jensen's inequality handles the posterior mean.
\end{proof}

\section{A range-preservation obstruction for full Gaussian path norm}
Minimizing the full Gaussian ReLU path norm can produce arbitrarily
large predictions between bounded observations. The following
construction concerns a prescribed finite sample; it makes no claim
about the probability of that sample geometry under an IID design.

For a univariate shallow ReLU implementation define
\[
 P(f)=\inf\left\{\sum_j|v_j|\sqrt{a_j^2+b_j^2}:
          f(x)=\sum_jv_j(a_jx+b_j)_+\ \text{on }[-1,1]\right\}.
\]
Every such representation satisfies
\[
 P(f)\ge\int_{(-1,1)}\sqrt{1+t^2}\,|f''|(dt), \tag{RG1}
\]
because an atom with knot $t=-b/a$ contributes second-derivative
mass $v|a|$, and
$\sqrt{1+t^2}|v a|=|v|\sqrt{a^2+b^2}$.
Cancellations only decrease the variation on the right.

Set $r=\sqrt3/2$ and $t_0=1/\sqrt3$. For
$0<\delta<1-r$, prescribe the four data pairs
\[
 (-r-\delta,-1),\quad(-r,0),\quad(r,0),\quad(r+\delta,-1).
\]
The three-neuron function
\[
 F_\delta(x)=\delta^{-1}
 \{r-t_0-(-x-t_0)_+-(x-t_0)_+\}
\]
interpolates them and has path norm at most
\[
 P(F_\delta)\le (r+\sqrt3)/\delta. \tag{RG2}
\]
Its middle plateau has height $(r-t_0)/\delta$.

More strongly, \emph{every} minimum-path-norm interpolant has
supremum magnitude at least $c/\delta$ for an absolute $c>0$.
To prove this, choose one fixed sufficiently small $\ell>0$ such that
\[
 a=\sqrt{1+(r-\ell)^2},\qquad 2a>r+\sqrt3.
\]
This is possible because
$2\sqrt{1+r^2}=\sqrt7>3\sqrt3/2=r+\sqrt3$.
Let $\psi_L$ be the triangular hat of height $a$, supported on
$[-r-\delta,-r+\ell]$, peaking at $-r$; let $\psi_R$ be its reflection.
Their supports are disjoint and
$0\le\psi_L+\psi_R\le\sqrt{1+x^2}$.
If an interpolant $f$ has $\|f\|_\infty\le M$, integration by parts gives
\[
 \int(\psi_L+\psi_R)\,df''
 =a\left[-\frac2\delta+
   \frac{f(-r+\ell)+f(r-\ell)}{\ell}\right].
\]
Consequently (RG1) implies
\[
 P(f)\ge \frac{2a}{\delta}-\frac{2aM}{\ell}. \tag{RG3}
\]
A minimum-path-norm interpolant has norm at most (RG2), so (RG3) yields
\[
 M\ge
 \frac{\ell[\,2a-(r+\sqrt3)\,]}{2a}\,\delta^{-1}.
\]
A minimizer exists: normalized signed unit-ridge evaluation vectors
form a compact subset of $\mathbb R^4$, so bounded-mass measures
have compact prediction sets. An optimum can be reduced to at most
four atoms by a linear-dependence argument.
Thus the statement concerns actual minimum-norm interpolants, not
only a particular nonminimal construction.

The example has a macroscopic unsampled central interval and two
very close pairs near its ends. Under a well-spread large IID sample
that geometry is rare. Moreover a likelihood with appreciable
regularization may decline to fit the close-pair label changes.
A successful posterior-specific improvement must quantify one of
those facts; a deterministic range-preservation assertion alone
cannot provide it.

\begin{corollary}[The actual posterior mean can have unbounded range]
Use the four observations in the preceding construction, so $n=4$.
There are absolute constants $c,c',\delta_0>0$ such that, for
$0<\delta<\delta_0$, $\lambda=\delta^2$, every $N\ge4$, and every
$0<\nu\le1$, the ordinary Gaussian posterior mean satisfies
\[
 \max\{\bar z(-r+\ell),\bar z(r-\ell)\}\ge c/\delta. \tag{RG4}
\]
Here $r$ and $\ell$ are the fixed constants in the preceding proof.
Moreover, against the uniform test-input law on $[-1,1]$ and any test
labels bounded by one, its squared test risk is at least
$c'/\delta^2$. The posterior-average squared test risk has the same
lower bound.

This is a deterministic-training-sample statement. It is not a lower
bound with a stated probability under an IID sample from that test law,
and it is not an impossibility theorem for other likelihood schedules.
\end{corollary}
\begin{proof}
Write $K=r+\sqrt3$. The explicit interpolant has zero empirical loss
and path variation at most $K/\delta$. Proposition (PR2) gives
\[
 \mathbb E_\Pi\widehat R(z)+2\lambda\mathbb E_\Pi P
 \le2\lambda K/\delta+C\lambda\nu
             \log\frac e{\lambda\nu}.
\]
For $\lambda=\delta^2$ and $0<\nu\le1$, this implies, uniformly in
both width and scale,
\[
 \mathbb E_\Pi\widehat R(z)\le C\delta,
 \qquad
 \mathbb E_\Pi P\le K/\delta+C\log(e/\delta). \tag{RG5}
\]
Let $f_i=z(X_i)$ in the order of the four training observations, and
let $f_L=z(-r+\ell)$, $f_R=z(r-\ell)$. The triangular-hat identity,
without assuming exact interpolation, is
\[
 \frac1a\int(\psi_L+\psi_R)\,dz''
 =\frac{f_1+f_4-f_2-f_3}{\delta}
   +\frac{f_L+f_R-f_2-f_3}{\ell}.
\]
Since the negative of this integral is at most $P$, rearrangement
and posterior expectation give
\begin{align*}
 \bar z(-r+\ell)+\bar z(r-\ell)
 &\ge \frac{2\ell}{\delta}
       -\frac{\ell}{a}\mathbb E_\Pi P
       -C(1+\delta^{-1})
                       \sqrt{\mathbb E_\Pi\widehat R(z)}\\
 &\ge\frac{\ell(2-K/a)}{\delta}
       -C\delta^{-1/2}-C\log(e/\delta).
\end{align*}
The leading coefficient is strictly positive by the choice of
$\ell$. For sufficiently small $\delta$, this proves (RG4).
All expectations are finite because the Gaussian posterior density
is bounded by $Z^{-1}$ times its prior density.

Every normalized ridge is $1$-Lipschitz, so the posterior mean is
Lipschitz with constant at most $\mathbb E_\Pi P\le2K/\delta$
for small $\delta$. A value at least $c/\delta$ at either fixed
interior point therefore persists at height $c/(2\delta)$ on an
interval of fixed positive length. For bounded test labels and
small enough $\delta$, integrating the squared error on that interval
gives $c'/\delta^2$. Jensen's inequality gives the posterior-average
risk lower bound as well.
\end{proof}

\section{Conditional readout uncertainty inside sample gaps}
The following bound separates uncertainty in the last Gaussian layer
from uncertainty in the hidden features. It applies to the actual
posterior and helps locate what remains missing from a sharper
population argument.

\begin{proposition}[Gaussian readout variance between observations]
Condition on the hidden weights $W=((a_j,b_j))_{j=1}^N$ and the sample.
Let $X_-<X_+$ be two observed inputs and let $x\in[X_-,X_+]$.
Write $h=X_+-X_-$ and
$Q_W=N^{-1}\sum_j(a_j^2+b_j^2)$.
Then
\[
 \operatorname{Var}_{\Pi}(z(x)\mid W,S)
 \le \frac{2n\lambda\nu}{N}+
       \frac{2\nu Q_Wh^2}{N}. \tag{RV1}
\]
For bounded labels, averaging over the hidden-weight posterior gives
\[
 \mathbb E_{W\mid S}\operatorname{Var}_{\Pi}(z(x)\mid W,S)
 \le\frac{2n\lambda\nu}{N}
 +\frac{2\nu h^2}{N}\left(3\nu+\frac1{2\lambda}\right). \tag{RV2}
\]
The same statement holds with $1$ replaced by $\widehat{Y^2}$ in
the last numerator when the observed labels are not bounded.
\end{proposition}
\begin{proof}
Let $H$ be the training feature matrix and let
$\sigma^2=n\lambda\nu/N$ be the actual observation variance.
The conditional Gaussian readout posterior has covariance
\[
 \Sigma=(N^2\nu^{-1}I+\sigma^{-2}H^\top H)^{-1}.
\]
Consequently $\Sigma\preceq(\nu/N^2)I$ and
$H\Sigma H^\top\preceq\sigma^2 I_n$.
Write $t=(x-X_-)/h$ and
\[
 h_W(x)=(1-t)h_W(X_-)+t h_W(X_+)+e(x).
\]
The ReLU Lipschitz property gives
$|e_j(x)|\le |a_j|h$.
The conditional variance of the first two terms is at most
$\sigma^2((1-t)^2+t^2)\le\sigma^2$, while that of the residual
is at most $(\nu/N^2)h^2\sum_j a_j^2\le\nu Q_Wh^2/N$.
Using $\operatorname{Var}(U+V)\le2\operatorname{Var}U+
2\operatorname{Var}V$ proves (RV1).

For completeness, Gaussian radial integration by parts gives
\[
 \mathbb E_\Pi Q_N
 =3\nu-\frac2\lambda\mathbb E_\Pi P_n[z(z-Y)]
 \le3\nu+\frac{\widehat{Y^2}}{2\lambda}.
\]
Indeed the Gaussian density is proportional to
$e^{-\omega Q_N/2}$, the parameter dimension is $3N$, and scaling
all parameters scales $z$ quadratically. Since $Q_W\le Q_N$,
averaging (RV1) proves (RV2).
\end{proof}

For a design with maximum adjacent sample gap
$h_{\max}=O(\log(n/\delta)/n)$ and $N\ge n$, this conditional
variance is at most
$O(\lambda+\log^2(n/\delta)/(n^3\lambda))$ in the sample range.
This does not control the variation of the conditional posterior
mean as the hidden weights vary. It also makes no assertion about
extrapolation beyond the smallest and largest observed inputs.
Consequently (RV2) alone does not bound the full posterior's
population squared error.


\section{What the reductions leave open}
The training-vector result applies at every finite width $N\ge n$,
with constants independent of width. In the strong-prior limit
$\nu\to0$ at fixed $\lambda$, both its objective gap and its
training-prediction discrepancy vanish. At fixed $\nu$, they also
vanish as $\lambda\to0$. These are statements about empirical
prediction and penalized empirical loss; the fitted vector and its
minimum-mass function extensions need not have stable extrapolation.

The four-point construction rules out a deterministic claim that
bounded labels force a bounded minimum-path fit or bounded posterior
mean. Its proof uses a large empty interval together with nearby
contradictory observations. A sharper random-design analysis could
still exploit the rarity of that configuration, or the regularization
cost of fitting close contradictory labels. The conditional variance
calculation shows that one must also distinguish Gaussian readout
uncertainty from the variation of the readout's conditional mean as
hidden weights change. No new population upper rate is claimed here.

\end{document}
